\chapter{DISCRETE-TIME MODIFIED STATE OBSERVER}
\label{ch:dmso}

The Discrete-Time Modified State Observer (DMSO) is an extension of the Modified
State Observer of Rajagopal et al.\ \cite{rajagopal2009mso} to discrete time.
This chapter states the problem, gives the observer and its adaptation law, and
establishes uniform ultimate boundedness of the state estimation error and the
functional approximation error by Lyapunov analysis.

\rev{The derivation evaluates the Lyapunov first difference exactly before any
bounding is performed. This is not a stylistic preference. The adaptation law is
constructed so that the cross terms coupling the weight error to the state
estimation error and to the network reconstruction error cancel identically, and
a term-by-term bound applied in advance destroys those cancellations. The
consequences are recorded in Appendix~\ref{app:errata}.}

\section{NOTATION}

$\norm{\cdot}$ denotes the Euclidean norm on vectors and the induced $2$-norm on
matrices; $\norm{\cdot}_F$ denotes the Frobenius norm, $\norm{M}_F^2=\Tr(M^TM)$.
For square $M$, $\sigmax(M)$ is the largest singular value and $\rho(M)$ the
spectral radius. Recall $\rho(M)\le\sigmax(M)=\norm{M}$, with equality only when
$M$ is normal \cite{horn2013matrix}. $\lmin(\cdot)$ and $\lmax(\cdot)$ denote
extreme eigenvalues of a symmetric matrix. Properties of semi-orthogonal
matrices and of the Frobenius norm follow Abadir and Magnus
\cite{abadir2005matrix}.

\section{PROBLEM STATEMENT}
\label{sec:problem}

\begin{assumption}[Plant]\label{as:plant}
The plant evolves as
\begin{equation}\label{eq:plant}
x_{k+1}=Fx_k+Gu_k+\Bm f(x_k),\qquad y_k=x_k,
\end{equation}
with $x_k\in\R^n$ the state, $u_k\in\R^m$ the control, $F\in\R^{n\times n}$ the
discretized system matrix, $G\in\R^{n\times m}$ the discretized input matrix,
and $f:\R^n\to\R^p$ an unknown smooth vector of uncertainties, $p$ being the
number of states carrying uncertainty.
\end{assumption}

\begin{assumption}[Selection matrix]\label{as:B}
$\Bm\in\R^{n\times p}$ contains exactly one unit entry per column, all other
entries zero, with no two columns equal. Its columns are therefore distinct
standard basis vectors of $\R^n$, so
\begin{equation}\label{eq:Bprops}
\Bm^T\Bm=I_p,\qquad \norm{\Bm}=1,\qquad \Bm\Bm^T=\Pi,
\end{equation}
where $\Pi$ is the orthogonal projector onto the uncertain coordinates. Only the
left identity in \eqref{eq:Bprops} is used below; note that
$\Bm\Bm^T\neq I_n$ whenever $p<n$.
\end{assumption}

\begin{assumption}[Functional approximation]\label{as:nn}
On the compact operating set there exist ideal weights $W\in\R^{l\times p}$ and
a basis $\phi:\R^n\to\R^l$ with
\begin{equation}\label{eq:fla}
f(x_k)=W^T\phi(x_k)+\epsilon_k,\qquad
\norm{\epsilon_k}\le\epsilon_N,\qquad
\norm{\phi(x_k)}\le\phi_{\max},
\end{equation}
$\epsilon_N$ and $\phi_{\max}$ known and finite. Write $\phi_k:=\phi(x_k)$. The
universal approximation property of two-layer networks is due to Hornik et al.\
\cite{hornik1989multilayer}; Sadegh \cite{sadegh1993nonlinear} extended it to the
single-layer functional link form used here, in which one layer suffices provided
$\phi$ is selected as a basis.
\end{assumption}

\begin{assumption}[Bounded ideal weights]\label{as:WM}
$\norm{W}_F\le W_M<\infty$. Used only in Section~\ref{sec:sigmod}.
\end{assumption}

\section{OBSERVER AND ADAPTATION LAW}

The observer is the modified Luenberger form
\begin{equation}\label{eq:obs}
\xhat_{k+1}=F\xhat_k+Gu_k+\Bm\fhat(x_k)+\Km\bigl(y_k-\xhat_k\bigr),
\qquad
\fhat(x_k)=\What_k^T\phi_k,
\end{equation}
with $\Km\in\R^{n\times n}$ a constant design gain and $\What_k\in\R^{l\times p}$
the weight estimate. Define the error-injection matrix
\begin{equation}\label{eq:A}
A:=F-\Km .
\end{equation}
Because $y_k=x_k$, $\Km$ is entirely at the designer's disposal and $A$ may be
assigned arbitrarily, including $A=0$; see Corollary~\ref{cor:deadbeat}.

The adaptation law is
\begin{equation}\label{eq:update}
\What_{k+1}=\What_k+\Gamma\,\phi_k\,e_{k+1}^T\Bm,\qquad \Gamma>0,
\end{equation}
where $e_k:=x_k-\xhat_k$ and $\Wt_k:=W-\What_k$. The law is causal: $\What_{k+1}$
is formed at step $k+1$, at which time $e_{k+1}$ is available.

Three quantities recur and are named once:
\begin{equation}\label{eq:shorthand}
\psi_k:=\Wt_k^T\phi_k\in\R^p,\qquad
\zeta_k:=\psi_k+\epsilon_k\in\R^p,\qquad
a_k:=\Bm^TAe_k\in\R^p .
\end{equation}
Here $\psi_k$ is the \emph{functional} approximation error, the quantity the
network must drive small, and $\zeta_k$ the total uncertainty reconstruction
error. Since $\norm{\Bm}=1$, $\norm{a_k}\le\norm{Ae_k}$.

\section{ERROR DYNAMICS}

\begin{lemma}[State estimation error]\label{lem:err}
Under Assumptions~\ref{as:plant}--\ref{as:nn},
\begin{equation}\label{eq:errdyn}
e_{k+1}=Ae_k+\Bm\zeta_k .
\end{equation}
\end{lemma}

\begin{proof}
Subtracting \eqref{eq:obs} from \eqref{eq:plant} with $y_k=x_k$,
\[
e_{k+1}=Fe_k-\Km e_k+\Bm\bigl(f(x_k)-\fhat(x_k)\bigr)
       =Ae_k+\Bm\bigl(W^T\phi_k+\epsilon_k-\What_k^T\phi_k\bigr),
\]
and $W^T\phi_k-\What_k^T\phi_k=\Wt_k^T\phi_k=\psi_k$.
\end{proof}

\begin{lemma}[Weight error]\label{lem:werr}
With $\xi_k:=a_k+\zeta_k$,
\begin{equation}\label{eq:wdyn}
\Wt_{k+1}=\Wt_k-\Gamma\phi_k\xi_k^T
=\bigl(I-\Gamma\phi_k\phi_k^T\bigr)\Wt_k-\Gamma\phi_k\bigl(a_k+\epsilon_k\bigr)^T .
\end{equation}
\end{lemma}

\begin{proof}
From \eqref{eq:update}, $\Wt_{k+1}=\Wt_k-\Gamma\phi_k e_{k+1}^T\Bm$. By
Lemma~\ref{lem:err} and $\Bm^T\Bm=I$,
\[
e_{k+1}^T\Bm=\bigl(Ae_k+\Bm\zeta_k\bigr)^T\Bm
=e_k^TA^T\Bm+\zeta_k^T=a_k^T+\zeta_k^T=\xi_k^T .
\]
The second equality follows from
$\Gamma\phi_k\zeta_k^T=\Gamma\phi_k\phi_k^T\Wt_k+\Gamma\phi_k\epsilon_k^T$.
\end{proof}

\section{LYAPUNOV STABILITY PROOF}
\label{sec:proof}

Take the Lyapunov candidate
\begin{equation}\label{eq:V}
V_k=e_k^Te_k+\frac{1}{\Gamma}\Tr\bigl(\Wt_k^T\Wt_k\bigr)
   =\norm{e_k}^2+\frac{1}{\Gamma}\norm{\Wt_k}_F^2\;\ge\;0 .
\end{equation}
\rev{The adaptation gain $\Gamma$ itself serves as the weighting. No independent
Lyapunov constant is introduced; introducing one and then tying it to $\Gamma$
by a later choice is what places the argument of revision~7 on the boundary of
its own admissibility condition.}

\subsection{The Exact First Difference}

\begin{proposition}[Exact difference]\label{prop:exact}
Let $\gamma_k:=\Gamma\norm{\phi_k}^2$. Then
\begin{equation}\label{eq:exact}
\Delta V_k=-\norm{e_k}^2
+\norm{Ae_k}^2+2a_k^T\epsilon_k+\norm{\epsilon_k}^2
-\norm{\psi_k}^2
+\gamma_k\norm{a_k+\zeta_k}^2 .
\end{equation}
\end{proposition}

\begin{proof}
\emph{Weight term.} From Lemma~\ref{lem:werr},
\[
\Tr\bigl(\Wt_{k+1}^T\Wt_{k+1}\bigr)
=\Tr\bigl(\Wt_k^T\Wt_k\bigr)
-2\Gamma\Tr\bigl(\Wt_k^T\phi_k\xi_k^T\bigr)
+\Gamma^2\Tr\bigl(\xi_k\phi_k^T\phi_k\xi_k^T\bigr).
\]
Since $\Tr(\Wt^T\phi\xi^T)=\xi^T\Wt^T\phi=\psi^T\xi$ and
$\Tr(\xi\phi^T\phi\xi^T)=\norm{\phi}^2\norm{\xi}^2$,
\begin{equation}\label{eq:wdiff}
\frac{1}{\Gamma}\Bigl[\norm{\Wt_{k+1}}_F^2-\norm{\Wt_k}_F^2\Bigr]
=-2\psi_k^T\xi_k+\gamma_k\norm{\xi_k}^2 .
\end{equation}

\emph{State term.} From Lemma~\ref{lem:err} and $\Bm^T\Bm=I$,
\begin{equation}\label{eq:ediff}
\norm{e_{k+1}}^2-\norm{e_k}^2
=\norm{Ae_k}^2-\norm{e_k}^2+2a_k^T\zeta_k+\norm{\zeta_k}^2,
\end{equation}
using $2e_k^TA^T\Bm\zeta_k=2a_k^T\zeta_k$.

\emph{Combine.} Adding \eqref{eq:wdiff} and \eqref{eq:ediff} with
$\xi_k=a_k+\zeta_k$,
\[
\Delta V_k=\norm{Ae_k}^2-\norm{e_k}^2
+\underbrace{2a_k^T\zeta_k-2\psi_k^Ta_k}_{\text{(i)}}
+\underbrace{\norm{\zeta_k}^2-2\psi_k^T\zeta_k}_{\text{(ii)}}
+\gamma_k\norm{a_k+\zeta_k}^2 .
\]
For (i), $2a_k^T(\zeta_k-\psi_k)=2a_k^T\epsilon_k$. For (ii), with
$\zeta_k=\psi_k+\epsilon_k$,
\[
\norm{\psi_k+\epsilon_k}^2-2\psi_k^T(\psi_k+\epsilon_k)
=\norm{\psi_k}^2+2\psi_k^T\epsilon_k+\norm{\epsilon_k}^2
-2\norm{\psi_k}^2-2\psi_k^T\epsilon_k
=-\norm{\psi_k}^2+\norm{\epsilon_k}^2 .
\]
Substituting gives \eqref{eq:exact}.
\end{proof}

The mechanism is now visible. Every cross term between $\psi_k$ and
$\epsilon_k$, and every cross term between $\psi_k$ and $a_k$, cancels exactly.
The negative term $-\norm{\psi_k}^2$ is produced by the adaptation law and is
not an artifact of any inequality. The only price paid for adaptation is the
single term $\gamma_k\norm{a_k+\zeta_k}^2$, which is small when $\Gamma$ is
small. Small adaptation gains must therefore be admissible on structural
grounds, and a condition that excludes them indicates an error upstream.

\subsection{Main Result}

\begin{theorem}[Uniform ultimate boundedness]\label{thm:main}
Let Assumptions~\ref{as:plant}--\ref{as:nn} hold. Set
$\bar\sigma:=\sigmax(A)$ and $\bar\gamma:=\Gamma\phi_{\max}^2$. Fix $\theta>0$
and define
\begin{equation}\label{eq:cs}
c_1:=1+\bar\gamma\Bigl(1+\tfrac{1}{\theta}\Bigr),\qquad
c_2:=1-\bar\gamma\bigl(1+\theta\bigr),\qquad
\rho:=\bar\sigma\sqrt{c_1}.
\end{equation}
If
\begin{align}
&\text{\rm(C1)}\quad c_2>0,
&&\text{i.e.}\quad \Gamma\phi_{\max}^2<\frac{1}{1+\theta},
\label{eq:C1}\\
&\text{\rm(C2)}\quad \rho<1,
&&\text{i.e.}\quad
\sigmax(A)^2\Bigl[1+\Gamma\phi_{\max}^2\bigl(1+\tfrac{1}{\theta}\bigr)\Bigr]<1,
\label{eq:C2}
\end{align}
then
\begin{equation}\label{eq:master}
\Delta V_k\;\le\;-\norm{e_k}^2
+c_1\bigl(\bar\sigma\norm{e_k}+\epsilon_N\bigr)^2
-c_2\norm{\psi_k}^2,
\end{equation}
and $\Delta V_k<0$ whenever
\begin{equation}\label{eq:bounds}
\norm{e_k}>b_e:=\frac{\sqrt{c_1}\,\epsilon_N}{1-\rho}
\qquad\text{or}\qquad
\norm{\psi_k}>b_\psi:=\epsilon_N\sqrt{\frac{c_1}{c_2\bigl(1-\rho^2\bigr)}} .
\end{equation}
The state estimation error $e_k$ and the functional approximation error
$\psi_k=\Wt_k^T\phi_k$ are uniformly ultimately bounded with ultimate bounds
$b_e$ and $b_\psi$.
\end{theorem}

\begin{proof}
Begin from \eqref{eq:exact}. Since $\norm{\Bm}=1$, $\norm{a_k}\le\norm{Ae_k}$, so
\[
\norm{Ae_k}^2+2a_k^T\epsilon_k+\norm{\epsilon_k}^2
\le\bigl(\norm{Ae_k}+\epsilon_N\bigr)^2 .
\]
Write $s_k:=\norm{Ae_k}+\epsilon_N$ and $w_k:=\norm{\psi_k}$. Then
$\norm{a_k+\zeta_k}\le\norm{Ae_k}+\norm{\psi_k}+\norm{\epsilon_k}\le s_k+w_k$,
and $\gamma_k\le\bar\gamma$, so
\[
\Delta V_k\le-\norm{e_k}^2+s_k^2-w_k^2+\bar\gamma(s_k+w_k)^2 .
\]
Applying Young's inequality \cite{young1912classes} once, in the form
$(s+w)^2\le(1+\tfrac{1}{\theta})s^2+(1+\theta)w^2$,
\[
\Delta V_k\le-\norm{e_k}^2+c_1s_k^2-c_2w_k^2 ,
\]
and $\norm{Ae_k}\le\bar\sigma\norm{e_k}$ gives \eqref{eq:master}.

For the first branch of \eqref{eq:bounds}, $c_2>0$ by (C1), so it suffices that
$c_1(\bar\sigma\norm{e_k}+\epsilon_N)^2<\norm{e_k}^2$, that is,
$\sqrt{c_1}(\bar\sigma\norm{e_k}+\epsilon_N)<\norm{e_k}$, that is,
$\norm{e_k}(1-\rho)>\sqrt{c_1}\epsilon_N$. This requires $\rho<1$, which is (C2).

For the second branch, let
$g(t):=c_1(\bar\sigma t+\epsilon_N)^2-t^2
=(\rho^2-1)t^2+2c_1\bar\sigma\epsilon_N t+c_1\epsilon_N^2$
for $t=\norm{e_k}\ge0$. By (C2) the leading coefficient is negative, so $g$ is
maximized at $t^\star=c_1\bar\sigma\epsilon_N/(1-\rho^2)$ with
\[
g(t^\star)=c_1\epsilon_N^2+\frac{c_1^2\bar\sigma^2\epsilon_N^2}{1-\rho^2}
=c_1\epsilon_N^2\Bigl[1+\frac{\rho^2}{1-\rho^2}\Bigr]
=\frac{c_1\epsilon_N^2}{1-\rho^2}.
\]
Hence $\Delta V_k\le g(\norm{e_k})-c_2w_k^2<0$ whenever
$c_2w_k^2>c_1\epsilon_N^2/(1-\rho^2)$, which is $\norm{\psi_k}>b_\psi$.

Since $V_k\ge0$ and $\Delta V_k<0$ outside
$\{\norm{e_k}\le b_e\}\cap\{\norm{\psi_k}\le b_\psi\}$, the standard Lyapunov
extension theorem for discrete-time systems
\cite{lewis1999neural,jagannathan2006neural} gives uniform ultimate
boundedness; see also Khalil \cite{khalil2002nonlinear} for the continuous-time
counterpart.
\end{proof}

\begin{corollary}[Simple sufficient conditions]\label{cor:theta1}
Take $\theta=1$, so $c_1=1+2\bar\gamma$ and $c_2=1-2\bar\gamma$. Then
Theorem~\ref{thm:main} applies provided
\begin{equation}\label{eq:simple}
\sigmax(A)<1
\quad\text{and}\quad
0<\Gamma<\frac{1}{\phi_{\max}^2}
\min\left\{\frac{1}{2},\;
\frac{1-\sigmax(A)^2}{2\,\sigmax(A)^2}\right\}.
\end{equation}
As $\Gamma\to0^+$ the ultimate bounds tend to
\begin{equation}\label{eq:limbounds}
b_e\to\frac{\epsilon_N}{1-\sigmax(A)},\qquad
b_\psi\to\frac{\epsilon_N}{\sqrt{1-\sigmax(A)^2}} .
\end{equation}
\end{corollary}

\begin{proof}
With $\theta=1$, (C1) is $\bar\gamma<1/2$ and (C2) is
$\bar\sigma^2(1+2\bar\gamma)<1$, i.e.\
$\bar\gamma<(1-\bar\sigma^2)/(2\bar\sigma^2)$, feasible for some
$\bar\gamma>0$ if and only if $\bar\sigma<1$. Substituting
$\bar\gamma=\Gamma\phi_{\max}^2$ gives \eqref{eq:simple}; the limits follow by
setting $c_1,c_2\to1$ and $\rho\to\bar\sigma$ in \eqref{eq:bounds}.
\end{proof}

\begin{corollary}[Deadbeat error injection]\label{cor:deadbeat}
Since $y_k=x_k$, choosing $\Km=F$ gives $A=0$. Working directly from
\eqref{eq:exact} with $a_k=0$,
\[
\Delta V_k=-\norm{e_k}^2-\bigl(1-\gamma_k\bigr)\norm{\psi_k}^2
+2\gamma_k\psi_k^T\epsilon_k+\bigl(1+\gamma_k\bigr)\norm{\epsilon_k}^2 ,
\]
so for any $\bar\gamma<1$, completing the square on the $\psi_k$ terms gives
\[
\Delta V_k\le-\norm{e_k}^2+\frac{\epsilon_N^2}{1-\bar\gamma},
\]
and $\norm{e_k}$ is ultimately bounded by $\epsilon_N/\sqrt{1-\bar\gamma}$. No
condition on $A$ is required and the restriction on $\Gamma$ relaxes to
$\bar\gamma<1$.
\end{corollary}

\begin{remark}[Interpretation]
Condition (C2) requires the error-injection matrix to be a contraction in the
$2$-norm; (C1) requires the adaptation gain to be small relative to the basis
magnitude. Faster error injection buys tolerance for a larger learning rate, and
a slower learning rate always helps. Both are what one would expect on physical
grounds. In the limit $\Gamma\to0^+$ the conditions collapse to
$\sigmax(A)<1$ and the bounds \eqref{eq:limbounds} reduce to the gain of the
error dynamics applied to the network reconstruction floor $\epsilon_N$, which
is the correct limiting answer.
\end{remark}

\begin{remark}[Non-normal $A$]\label{rem:nonnormal}
If $A$ is Schur but $\sigmax(A)\ge1$, Theorem~\ref{thm:main} does not apply in
the Euclidean norm. Since $\rho(A)<1$, for any $\delta>0$ there is a nonsingular
$T$ with $\norm{TAT^{-1}}\le\rho(A)+\delta<1$. Repeating the derivation with
$V_k=e_k^TT^TTe_k+\Gamma^{-1}\norm{\Wt_k}_F^2$ yields the same conclusions with
$\bar\sigma$ replaced by $\norm{TAT^{-1}}$, at the cost of a factor $\kappa(T)$
in the ultimate bounds. Since $\Km$ is fully free here, it is simpler in
practice to place $A$ so that $\sigmax(A)$ is directly small. Note that
$\lmax(A)$ is not an upper bound for $\norm{A}$ unless $A$ is normal, and may be
complex.
\end{remark}

\section{BOUNDEDNESS OF THE WEIGHT ESTIMATES}
\label{sec:weightbound}

Theorem~\ref{thm:main} bounds $\psi_k=\Wt_k^T\phi_k$, not $\Wt_k$. This is the
correct statement: the negative term in \eqref{eq:master} acts only along
$\phi_k$, so components of $\Wt_k$ orthogonal to the regressor direction are
unconstrained by \eqref{eq:master} and may drift. Two standard remedies apply.

\subsection{Persistency of Excitation}

\begin{lemma}[Sadegh \cite{sadegh1993nonlinear}]\label{lem:pe1}
Let $\Phi(k_1,k_0)$ denote the state-transition matrix of
$\Wt_{k+1}=M(k)\Wt_k$. If $\norm{\Phi(k_1,k_0)}<1$ for all
$k_1>k_0\ge0$, the homogeneous system is stable.
\end{lemma}

\begin{definition}[Persistency of excitation]\label{def:pe}
The sequence $\{\phi_k\}$ is persistently exciting if there exist $\beta>0$ and
an integer $L\ge1$ such that
\begin{equation}\label{eq:pe}
\lmin\left(\sum_{j=k}^{k+L-1}\phi_j\phi_j^T\right)\ge\beta,\qquad
\forall\,k\ge0 .
\end{equation}
\end{definition}

\begin{lemma}\label{lem:pe2}
If $M(k)=I-\Gamma\phi_k\phi_k^T$ with $0<\Gamma\norm{\phi_k}^2<2$ and
\eqref{eq:pe} holds, then the homogeneous part of \eqref{eq:wdyn} is uniformly
exponentially stable, and since the forcing term
$-\Gamma\phi_k(a_k+\epsilon_k)^T$ is bounded whenever $e_k$ is bounded, $\Wt_k$
is bounded.
\end{lemma}

\rev{The condition $0<\Gamma\norm{\phi_k}^2<2$ required here is exactly the
condition $\gamma_k<2$ that renders the $\psi_k$ coefficient in \eqref{eq:wdiff}
negative. In the present derivation it is not an independent assumption but a
consequence of the same inequality that drives Theorem~\ref{thm:main}, and
Corollary~\ref{cor:theta1} implies it.}

\subsection{$\sigma$-Modification}
\label{sec:sigmod}

Persistency of excitation is not always available, particularly during
regulation, where the regressor may be nearly constant. Replacing
\eqref{eq:update} by
\begin{equation}\label{eq:sigmod}
\What_{k+1}=\bigl(1-\Gamma\kappa\bigr)\What_k+\Gamma\phi_ke_{k+1}^T\Bm,
\qquad\kappa>0,
\end{equation}
gives $\Wt_{k+1}=\Wt_k-\Gamma\phi_k\xi_k^T+\Gamma\kappa\What_k$. The linear part
of the trace difference acquires
\begin{equation}\label{eq:sigmodterm}
2\kappa\Tr\bigl(\Wt_k^TW\bigr)-2\kappa\norm{\Wt_k}_F^2
\;\le\;-\kappa\norm{\Wt_k}_F^2+\kappa W_M^2 ,
\end{equation}
using Assumption~\ref{as:WM} and $2ab\le a^2+b^2$. This supplies negative
definiteness in $\Wt_k$ directly. The quadratic contributions of
$\Gamma\kappa\What_k$ are dominated for $\Gamma\kappa$ sufficiently small, and
the cost is a bias of order $\kappa W_M^2$ in the ultimate bound.
$\sigma$-modification \cite{ioannou1983adaptive,ioannou1984sigma} is the same
device used for the controller networks in Chapter~\ref{ch:control}; applying it to the observer network as well makes the
two stability arguments structurally consistent and is the recommended
implementation.

\section{OBSERVER OR FILTER}

With $y_k=x_k$ the full state is measured, so \eqref{eq:obs} is properly a
filter and uncertainty identifier rather than an observer in the
unmeasured-state sense. Nothing in the proof depends on this, but it should be
stated plainly, and it is the reason $\Km$ may be chosen freely in
Corollary~\ref{cor:deadbeat}. Extension to $y_k=Cx_k$ with $\operatorname{rank}
C<n$ requires $A=F-\Km C$ and a detectability assumption; the structure of
Proposition~\ref{prop:exact} carries over unchanged with
$a_k=\Bm^TAe_k$.
